\section{Exact Moore bounds and cycle exclusion}\label{sec:cycles}

We now formulate the two sides of this density--girth tension for the subgraph
induced by the vertices of degree \(\le4\).  An exact Moore bound converts its
edge excess into a short cycle, whereas the spectral certificates of
Section~\ref{sec:certificates} exclude every cycle in the same length range.

\subsection{Exact non-backtracking growth}

\begin{proposition}[Non-backtracking Moore inequality]\label{lem:nb-moore}
Let \(H\) be a finite simple graph with \(N\) vertices and \(\delta(H)\ge2\).
Write \(d_H(x)\) for the degree of \(x\) in \(H\), let
\(D_H=\sum_{x\in V(H)}d_H(x)\) be its degree sum, and let \(\ell\ge1\) be an
integer.  If \(H\) has no cycle of length \(\le2\ell\), then
\[
 D_H\sum_{i=0}^{\ell-1}(D_H-N)^iN^{\ell-1-i}\le N^\ell(N-1).
\]
\end{proposition}

\begin{proof}
Put \(\bar d=D_H/N\).  Since the girth of \(H\) is at least
\(2\ell+1\), the irregular Moore bound of Alon, Hoory, and
Linial~\cite{ahl2002} gives
\[
 N\ge 1+\bar d\sum_{i=0}^{\ell-1}(\bar d-1)^i.
\]
Here the odd-girth bound applies directly when the girth is odd; when it is
even, the corresponding even-girth bound is stronger.  Substituting
\(\bar d=D_H/N\), subtracting \(1\), and multiplying by \(N^\ell\) gives the
stated inequality.
\end{proof}

To apply Proposition~\ref{lem:nb-moore} to an induced subgraph with \(v\)
vertices and edge surplus at least \(t\), we use the numerical condition
\begin{equation}\label{eq:ball-condition}
 t(v-1)v^\ell<(v+t)\bigl((v+2t)^\ell-v^\ell\bigr).
\end{equation}
The next proposition shows that this condition survives passage to the
possibly smaller \(2\)-core and forces a short cycle.

\begin{proposition}[Induced-set Moore criterion]\label{prop:nb-induced-cycle}
Let \(G\) be a finite simple graph and let \(U\subseteq V(G)\) be nonempty.
Put \(v=|U|\), and suppose \(e(U)\ge v+t\) for an integer \(t\ge1\).  Let \(L\ge6\) and
\(\ell=\lfloor L/2\rfloor\).  If \eqref{eq:ball-condition} holds, then
\(G[U]\) contains a cycle of length \(\le L\).
\end{proposition}

\begin{proof}
Repeatedly delete vertices of degree \(\le1\) from \(G[U]\), and let
\(H\) be the resulting \(2\)-core.  Each deletion removes one vertex and
\(\le1\) edge, so \(e-|V|\) does not decrease.  Since initially
\(e(U)-v\ge t>0\), the process cannot delete every vertex.  Thus \(H\) is
nonempty and, if \(v'=|V(H)|\), then \(e(H)\ge v'+t\).  Its total degree
\(D\) therefore satisfies \(D\ge2(v'+t)\) and \(D-v'\ge v'+2t\).

After division by \(tv^\ell\), condition \eqref{eq:ball-condition} is
equivalent to
\[
 v-1<2\left(1+\frac tv\right)
 \sum_{i=0}^{\ell-1}\left(1+\frac{2t}{v}\right)^i.
\]
The left-hand side decreases and the right-hand side increases as \(v\)
decreases, because both \(t/v\) and \(2t/v\) increase.  Since \(v'\le v\),
the same strict inequality holds with \(v'\) in place of \(v\).

Suppose that \(H\) has no cycle of length \(\le L\).  Since
\(2\ell\le L\), Proposition~\ref{lem:nb-moore} and the preceding degree
bounds give, term by term,
\[
 2(v'+t)\sum_{i=0}^{\ell-1}(v'+2t)^i(v')^{\ell-1-i}
 \le (v')^\ell(v'-1).
\]
Multiplying by \(t\) and using
\[
 2t\sum_{i=0}^{\ell-1}(v'+2t)^i(v')^{\ell-1-i}
 = (v'+2t)^\ell-(v')^\ell
\]
gives
\[
 (v'+t)\bigl((v'+2t)^\ell-(v')^\ell\bigr)
 \le t(v'-1)(v')^\ell,
\]
the reverse of \eqref{eq:ball-condition} for \(v'\), a contradiction.
\end{proof}

\subsection{Spectral cycle exclusion}

The Moore criterion produces a cycle of length \(\le L\) in \(G[S]\).  The
next proposition shows that the spectral certificates forbid such a cycle
whenever the census leaves enough boundary capacity.

\begin{proposition}[Short-cycle exclusion criterion]
\label{lem:cycle-exclusion}
Let \(G\) be a degree-\(3\)-separated obstruction, and let \(L\ge3\) be an
integer.  If \(12L+X\le2n\), then \(G[S]\) contains no cycle of length
\(\le L\).
\end{proposition}

\begin{proof}
Suppose that \(G[S]\) contains a cycle \(C\) of length \(k\le L\), and
put \(A=V(C)\), \(F=N(A)\setminus A\), and
\(B=V\setminus(A\cup F)\).  Write
\(D_C=\sum_{v\in A}d(v)\), \(f=|F|\), and \(b=|B|\).  The cycle uses two
incident edges at each of its \(k\) vertices.  Since every vertex of \(A\)
has degree \(\le4\), and every vertex of \(F\) has a neighbor in \(A\),
\begin{equation}\label{eq:cycle-slice}
  e(A,F)\le D_C-2k\le2k=2|A|,
  \qquad f\le D_C-2k.
\end{equation}

There are \(X+8\) degree-\(3\) vertices, by \eqref{eq:d3-count}, so the
number of degree-\(3\) vertices in \(B\) is \(\le X+8\).  Hence at least
\(b-X-8\) vertices of \(B\) have degree \(\ge4\), and
\(\sum_{v\in B}(d(v)-3)\ge b-X-8\).  Using
\eqref{eq:total-excess} and the exact excess \(D_C-3k\) of \(A\), we obtain
\[
\begin{aligned}
  e(B,F)
  &\le2f+\sum_{v\in F}(d(v)-3)\\
  &\le2f+(n-8)-(D_C-3k)-(b-X-8)\\
  &=4k-D_C+X+3f.
\end{aligned}
\]
Thus \(e(B,F)\le2b\) follows from \(5f+6k+X-D_C\le2n\).  By
\eqref{eq:cycle-slice}, its left-hand side is
\(\le4D_C-4k+X\le12k+X\le12L+X\le2n\).

Since \(k\le L\), the hypothesis also gives \(12k\le12L\le2n-X\le2n\),
so \(n\ge6k\).  As \(f\le2k\), we have
\(b\ge n-3k\ge3k>0\).  Thus \(A\) and \(B\) are nonempty; they are
nonadjacent by the definition of \(F\).  Equation~\eqref{eq:cycle-slice}
and the preceding calculation give their two required boundary bounds.
Proposition~\ref{lem:two-cluster} therefore gives \(\ac(G)\le2\),
contradicting the definition of an obstruction.
\end{proof}

Together with the edge surplus in Lemma~\ref{lem:S-density}, the two cycle
criteria reduce the proof for \(n\ge48\) to finding a length \(L\) for which
\eqref{eq:ball-condition} and \(12L+X\le2n\) hold simultaneously.

\section{Obstruction exclusion for \texorpdfstring{\(n\ge32\)}{n >= 32}}
\label{sec:obstruction-closure}

We now rule out the obstruction class in two overlapping ranges.  For
\(n\ge48\), the exact Moore criterion and the spectral cycle-exclusion
criterion are incompatible.  For \(32\le n\le49\), the incidence census
exceeds the available high-degree capacity.

\subsection{The Moore argument for \texorpdfstring{\(n\ge48\)}{n >= 48}}

We apply the two criteria from Section~\ref{sec:cycles} at the largest cycle
length permitted by the spectral budget.  The census will place the resulting
parameters in the region
\(43\le t\le v\) and \(5v+41\le50\ell+3t\).  We first isolate the power
estimate needed when \(4\le\ell\le7\): the case \(\ell=4\) is treated
directly, while \(5\le\ell\le7\) uses the first three nonconstant binomial
terms.  For \(\ell\ge8\), the fourth term alone will be strong enough.

\begin{lemma}[Small-exponent power growth]\label{lem:small-exponent-power}
Let \(t,v,\ell\) be integers satisfying
\[
  43\le t\le v\le107,
  \qquad 4\le\ell\le7,
  \qquad 5v+41\le50\ell+3t.
\]
Then
\begin{equation}\label{eq:small-exponent-power}
 (t-9)v^\ell\le(v+2t)^\ell.
\end{equation}
\end{lemma}

\begin{proof}
Put \(M=50\ell+3t-41\).  The linear hypothesis gives \(5v\le M\), and
\(t\le v\) gives \(5t\le M\).

First suppose that \(\ell=4\), so \(M=159+3t\).  The inequality
\(5t\le M\) gives \(t\le79\).  To verify
\((t-9)M^4\le(M+10t)^4\), set \(u=t-43\), so \(0\le u\le36\).  Direct
expansion gives
\begin{align*}
 &(M+10t)^4-(t-9)M^4\\
 &\quad={}
 u^3\bigl(773272-81u^2-5297u\bigr)
 +83801688u^2+2621645152u+31854951952.
\end{align*}
The coefficient in parentheses decreases on \([0,36]\) and has value
\(477604\) at \(u=36\), so the expression is nonnegative.  Since
\(M,v>0\) and \(5v\le M\),
\[
 t-9\le\left(\frac{M+10t}{M}\right)^4
 \le\left(\frac{5v+10t}{5v}\right)^4
 =\left(\frac{v+2t}{v}\right)^4,
\]
which proves \eqref{eq:small-exponent-power} for \(\ell=4\).

Now suppose that \(5\le\ell\le7\).  We first prove
\begin{align}
 3(t-10)v^3\le{}&6\ell t v^2
 +6\ell(\ell-1)t^2v
 +4\ell(\ell-1)(\ell-2)t^3. \label{eq:small-exponent-cubic}
\end{align}
For real \(x\), define
\begin{align*}
 \Phi(x)={}&30\ell t x^2+150\ell(\ell-1)t^2x
       +500\ell(\ell-1)(\ell-2)t^3\\
       &-3(t-10)x^3.
\end{align*}
The hypotheses give \(5t\le5v\le M\), and
\eqref{eq:small-exponent-cubic} is equivalent to \(\Phi(5v)\ge0\).  On
\([5t,M]\),
\[
 \Phi''(x)=60\ell t-18(t-10)x
 \le t(1320-90t)<0,
\]
so \(\Phi\) is concave and its minimum occurs at an endpoint.

Set \(a=\ell-5\) and \(b=t-43\).  Then \(0\le a\le2\), while
\(t\le v\le107\) gives \(0\le b\le64\).  At the left endpoint,
\[
 \frac{\Phi(5t)}{125t^3}
 =4a^3+54a^2+248a+54+3(79-b)\ge0.
\]
For the other endpoint, define
\begin{align*}
 Q_3(b)&=500b^3+72000b^2+3118500b+44471000,\\
 Q_2(b)&=6450b^3+872400b^2+36222750b+504355800,\\
 Q_1(b)&=23770b^3+3057300b^2+121507590b+1658324560.
\end{align*}
Then
\begin{align*}
 \Phi(M)={}&a^3Q_3(b)+a^2Q_2(b)+aQ_1(b)
 +b^3(10299-81b)\\
 &+2032188b^2+82837380b+1174137072.
\end{align*}
Every term is nonnegative because \(10299-81b\ge5115\).  Thus
\(\Phi(5v)\ge0\), proving \eqref{eq:small-exponent-cubic}.

Finally, the first three nonconstant terms of the binomial expansion give
\begin{align*}
3(v+2t)^\ell\ge{}&3v^\ell+6\ell t v^{\ell-1}
 +6\ell(\ell-1)t^2v^{\ell-2}\\
&+4\ell(\ell-1)(\ell-2)t^3v^{\ell-3}.
\end{align*}
After multiplying \eqref{eq:small-exponent-cubic} by
\(v^{\ell-3}\), the last three terms are at least
\(3(t-10)v^\ell\).  Hence
\(3(v+2t)^\ell\ge3(t-9)v^\ell\), proving the lemma.
\end{proof}

\begin{proposition}[Uniform Moore arithmetic]\label{prop:moore-arithmetic}
Suppose that \(n,X,h\) are nonnegative integers satisfying
\[
 n\ge48,
 \qquad h\le X,
 \qquad 10X+7h+186\le4n.
\]
Put
\[
 t=n-4-X-3h,\qquad v=n-h,\qquad
 L=\left\lfloor\frac{2n-X}{12}\right\rfloor,\qquad
 \ell=\left\lfloor\frac L2\right\rfloor.
\]
Then \(t\ge43\), \(L\ge8\), and \eqref{eq:ball-condition} holds.
\end{proposition}

\begin{proof}
The nested-floor identity
\(\lfloor\lfloor q/12\rfloor/2\rfloor=\lfloor q/24\rfloor\), applied to
\(q=2n-X\), gives \(\ell=\lfloor(2n-X)/24\rfloor\).  We may therefore write
\begin{equation}\label{eq:moore-remainder}
 2n-X=24\ell+r,\qquad 0\le r\le23.
\end{equation}
The census inequality, together with \(h\le X\), first gives
\[
 4t=4n-16-4X-12h\ge6X-5h+170\ge170,
\]
so \(t\ge43\); also \(v-t=4+X+2h\ge0\).  Moreover,
rearranging the same census inequality gives
\(10(2n-X)\ge16n+7h+186\ge954\).  Hence
\(2n-X\ge96\), \(\ell\ge4\), and \(L\ge8\).

Rewriting the census inequality using \eqref{eq:moore-remainder} gives
\begin{equation}\label{eq:moore-remainder-census}
 8X+7h+186\le48\ell+2r.
\end{equation}
The identity \(50\ell+3t-(5v+41)=26\ell-r-4X-4h-53\) shows that it remains
to prove \(4X+4h+r+53\le26\ell\).  Since \(h\le X\)
and \(r\le23\), \eqref{eq:moore-remainder-census} gives
\(15h\le48\ell-140\le15(4\ell-12)\), where the second inequality uses
\(\ell\ge4\).  Hence \(h\le4\ell-12\).
Using this and \(r\le23\) once more, we obtain
\[
\begin{aligned}
 2(4X+4h+r+53)
 &=8X+8h+2r+106\\
 &\le48\ell+h+4r-80\le52\ell.
\end{aligned}
\]
We have therefore established
\begin{equation}\label{eq:moore-region}
 43\le t\le v,\qquad \ell\ge4,\qquad
 5v+41\le50\ell+3t.
\end{equation}

Suppose first that \(\ell\le7\).  Equation
\eqref{eq:moore-remainder} gives \(2n-X\le191\), or
\(X\ge2n-191\), while the census gives \(10X\le4n-186\).  Therefore
\(20n-1910\le10X\le4n-186\), so \(n\le107\) and \(v\le107\).
Lemma~\ref{lem:small-exponent-power} now gives
\((t-9)v^\ell\le(v+2t)^\ell\).  Furthermore,
\[
 (t-10)(v+t)-\bigl(t(v-1)+1\bigr)
 =t^2-9t-10v-1>0,
\]
because its right-hand side is at least
\(43^2-9\cdot43-10\cdot107-1=391\).  It follows that
\[
\begin{aligned}
 t(v-1)v^\ell
 &<\bigl(t(v-1)+1\bigr)v^\ell\\
 &\le(v+t)(t-10)v^\ell\\
 &\le(v+t)\bigl((v+2t)^\ell-v^\ell\bigr),
\end{aligned}
\]
which is \eqref{eq:ball-condition}.

It remains to consider \(\ell\ge8\).  Put
\(c_\ell=\ell(\ell-1)(\ell-2)(\ell-3)\).  We claim that
\begin{equation}\label{eq:fourth-term-region}
 3v^4\le2c_\ell t^3.
\end{equation}
If \(t=43\), then \eqref{eq:moore-region} gives
\(v\le10\ell+17\).  Writing \(u=\ell-8\), direct expansion gives
\begin{align*}
2c_\ell\cdot43^3-3(10\ell+17)^4
={}&129014u^4+2970364u^3+22976314u^2\\
 &+59988164u+1555677>0.
\end{align*}
Together with \(v\le10\ell+17\), this proves
\eqref{eq:fourth-term-region} when \(t=43\).

Now suppose that \(t\ge44\), and set \(M_s=50\ell+3s-41\) for \(s\ge44\).
Equation \eqref{eq:moore-region} gives \(5v\le M_t\), while
\(t\le v\) gives \(M_t\ge5t\).  At \(t=44\), again writing
\(u=\ell-8\),
\begin{align*}
1250c_\ell\cdot44^3-3(50\ell+91)^4
={}&87730000u^4+2031980000u^3+15877835000u^2\\
&+42485217400u+4526254317>0.
\end{align*}
For fixed \(\ell\), the ratio \(M_s^4/s^3\) does not increase with
\(s\) while \(M_s\ge5s\).  Indeed, for \(y=1/s\),
\[
 (1+y)^3-\left(1+\frac{3y}{5}\right)^4
 =\frac{y}{625}\bigl(375+525y+85y^2-81y^3\bigr)\ge0,
\]
because \(0<y\le1\).  Since \(M_{s+1}=M_s+3\), this gives
\[
 \left(\frac{M_{s+1}}{M_s}\right)^4
 \le\left(1+\frac{3}{5s}\right)^4
 \le\left(1+\frac1s\right)^3.
\]
The difference \(M_s-5s=50\ell-41-2s\) decreases with \(s\), so
\(M_t\ge5t\) implies \(M_s\ge5s\) for every integer \(44\le s\le t\).
The preceding comparison says precisely that
\(M_{s+1}^4/(s+1)^3\le M_s^4/s^3\).  Thus the inequality at \(s=44\)
propagates to \(s=t\), giving
\(3M_t^4\le1250c_\ell t^3\).  Since \(5v\le M_t\),
\[
 625\cdot3v^4=3(5v)^4\le3M_t^4\le1250c_\ell t^3.
\]
Division by \(625\) proves \eqref{eq:fourth-term-region}.

The fourth nonconstant term of the binomial expansion now gives
\[
\begin{aligned}
3(v+2t)^\ell
&\ge3v^\ell+2c_\ell t^4v^{\ell-4}\\
&\ge3v^\ell+3t v^\ell
 =3(t+1)v^\ell.
\end{aligned}
\]
Consequently
\((v+2t)^\ell-v^\ell\ge t v^\ell\), and hence
\[
 t(v-1)v^\ell<(v+t)t v^\ell
 \le(v+t)\bigl((v+2t)^\ell-v^\ell\bigr).
\]
This is again \eqref{eq:ball-condition}.
\end{proof}

\begin{proposition}[No obstruction in the Moore range]
\label{prop:no-obstruction-ge48}
There is no degree-\(3\)-separated obstruction of order \(n\ge48\).
\end{proposition}

\begin{proof}
Suppose that \(G\) is such an obstruction.  Lemma~\ref{lem:moore-census}
and \eqref{eq:h-le-X} verify the hypotheses of
Proposition~\ref{prop:moore-arithmetic}.  Let
\[
  t=n-4-X-3h,\qquad
  v=n-h=|S|,\qquad
  L=\left\lfloor\frac{2n-X}{12}\right\rfloor,\qquad
  \ell=\left\lfloor\frac L2\right\rfloor.
\]
The proposition gives \(t\ge43\), \(L\ge8\), and
\eqref{eq:ball-condition}.  The set \(S\) is nonempty because it contains
the \(8+X\) vertices of degree \(3\), and Lemma~\ref{lem:S-density} gives
\(e(S)\ge|S|+t\).  Proposition~\ref{prop:nb-induced-cycle} therefore
produces a cycle in \(G[S]\) of length \(\le L\).  On the other hand, the
definition of \(L\) gives \(12L+X\le2n\), so
Proposition~\ref{lem:cycle-exclusion} excludes every such cycle.
This contradiction proves the proposition.
\end{proof}

\subsection{The incidence argument for
\texorpdfstring{\(32\le n\le49\)}{32 <= n <= 49}}

To cover the orders below \(48\), we return directly to the incidence census;
the resulting argument is valid throughout \(32\le n\le49\).  Retain
\(X,h,\tau(x),\tau_4\), and \(g\) from Section~\ref{sec:census}.  The vertices
counted by \(g\) are precisely those for which the star certificate does not
impose the bound \(\tau(x)\le d(x)-3\).

\begin{proposition}[Incidence-capacity contradiction]
\label{prop:incidence-capacity-contradiction}
There is no degree-\(3\)-separated obstruction of order
\(32\le n\le49\).
\end{proposition}

\begin{proof}
Suppose that \(G\) is such an obstruction.  We first derive the two capacity
bounds
\begin{align}
  24+2X&\le\tau_4+h+3g,              \label{eq:short-slots}\\
  7\tau_4+3X+32&\le4n.               \label{eq:short-choke}
\end{align}
If \(9d(x)\le n+15\), the star certificate gives
\(\tau(x)\le d(x)-3\); for a vertex counted by \(g\), the trivial bound
costs three additional units.  Thus
\(\sum_{d(x)\ge5}\tau(x)\le X+h+3g\).  Adding the degree-\(4\) contribution
and using \(\sum_{d(x)\ge4}\tau(x)=24+3X\) proves
\eqref{eq:short-slots}.

For \eqref{eq:short-choke}, the inequality \(36\le n+15\) and
Lemma~\ref{lem:star} show that every degree-\(4\) vertex has \(\le1\)
neighbor in \(D_3\).  Hence, for
\(O=\{x:d(x)=4,\ \tau(x)=1\}\), we have \(|O|=\tau_4\).
The set \(O\) is independent.  Indeed, adjacent members with distinct
degree-\(3\) neighbors contradict Lemma~\ref{lem:decorated-edge}(i), while
a shared neighbor creates a \((4,4,3)\)-triangle covered by
Lemma~\ref{lem:cycle-certificate}.

Each member of \(O\) sends its other three edges to vertices of degree
at least \(4\) outside \(O\).  By \eqref{eq:high-degree-sum}, the degree sum
over all vertices of degree at least \(4\) is \(4n-32-3X\); after subtracting
the \(4\tau_4\) contributed by \(O\), we obtain
\(3\tau_4\le4n-32-3X-4\tau_4\), which is \eqref{eq:short-choke}.

Multiply \eqref{eq:short-slots} by \(7\), then use
\eqref{eq:short-choke} and \eqref{eq:h-le-X}.  The result is
\begin{equation}\label{eq:short-aggregate}
  200+10X\le4n+21g.
\end{equation}

If \(39\le n\le49\), then \eqref{eq:exceptional-giants} gives
\(19g\le9X\), hence \(21g\le10X\).  Equation
\eqref{eq:short-aggregate} would force \(200\le4n\), contrary to
\(n\le49\).

Now let \(32\le n\le38\).  Three times
\eqref{eq:exceptional-hoarding}, together with
\eqref{eq:exceptional-giants}, gives
\[
 g(n-29)\le3(n-32)\le2(n-29),
\]
so \(g\le2\).  Equation \eqref{eq:short-aggregate} then gives
\(200\le4n+21g\le194\), again a contradiction.
\end{proof}

The two arguments overlap at orders \(48\) and \(49\) and together exclude
the obstruction class above order \(31\).

\begin{theorem}[Orders at least thirty-two]\label{thm:range-ge32}
If \(n\ge32\) and \(G\) is a graph on \(n\) vertices with
\(2(n-2)\) edges, then \(\ac(G)\le2\).
\end{theorem}

\begin{proof}
Suppose, for a contradiction, that \(\ac(G)>2\).  A vertex of degree
\(\le2\) would have a closed neighborhood of size \(\le3\), whose complement
is nonempty; Lemma~\ref{lem:low-degree} would then give \(\ac(G)\le2\).
Thus \(\delta(G)\ge3\).  Similarly, Lemma~\ref{lem:adjacent-three} shows that
no two degree-\(3\) vertices are adjacent.  Hence \(G\) is a
degree-\(3\)-separated obstruction.  If \(n\le49\), this contradicts
Proposition~\ref{prop:incidence-capacity-contradiction}; if \(n\ge50\), it
    contradicts Proposition~\ref{prop:no-obstruction-ge48}.
\end{proof}

Thus every graph of order \(n\ge32\) satisfies the desired bound.  Below this
threshold, the same degree count still forces a shared degree-\(4\) hub, but
additional local certificates are needed to exclude that configuration.
